\begin{tikzpicture}[
  scale=0.94,transform shape,
  node distance=5mm and 4mm,
  every node/.style={font=\sffamily\scriptsize,align=center},
  cat/.style={draw=uerjblue,rounded corners,fill=lightblue,
              text width=2.05cm,minimum height=0.82cm,inner sep=2pt}
]
  \node[cat] (cla) {material\\fragmentar};

  \node[cat,below left=7mm and -1mm of cla]
    (epi) {terrígeno};

  \node[cat,below right=7mm and -1mm of cla]
    (vul) {vulcanoclástico};

  \node[cat,below left=13mm and -1mm of vul]
    (pir) {piroclástico};

  \node[cat,below right=13mm and -1mm of vul]
    (ret) {retrabalhado\\ou ressedimentado};

  \draw[-{Latex[length=2mm]},uerjred,thick] (cla) -- (epi);
  \draw[-{Latex[length=2mm]},uerjred,thick] (cla) -- (vul);
  \draw[-{Latex[length=2mm]},uerjred,thick] (vul) -- (pir);
  \draw[-{Latex[length=2mm]},uerjred,thick] (vul) -- (ret);
\end{tikzpicture}